\documentclass{amsart}
\usepackage{amsmath,amssymb,amsthm,hyperref}

\newtheorem{theorem}{Theorem}
\newtheorem{proposition}{Proposition}
\newtheorem{lemma}{Lemma}
\newtheorem{corollary}{Corollary}
\theoremstyle{definition}
\newtheorem{definition}{Definition}
\newtheorem{assumption}{Assumption}
\theoremstyle{remark}
\newtheorem{remark}{Remark}

\begin{document}

\title{A Multilevel Mediation Model of Team Learning Behavior in Field Settings}
\maketitle

\begin{abstract}
We give an explicit, formally specified model of the antecedents of team
learning behavior in ongoing organizational work teams. The model
identifies factors at the individual, team, and organizational levels;
articulates the mechanism by which each acts; and singles out a shared
interpersonal-climate variable---team psychological safety, defined as the
shared belief that the team is safe for interpersonal risk taking---as the
mediator linking contextual supports (supportive leadership, contextual
support, coaching) to observed learning behavior. We state the model as a
recursive structural (path) system, derive the conditions under which the
climate variable is a genuine mediator (the indirect effect is nonzero and
the direct effect vanishes), and prove that, under the stated structural
assumptions, contextual supports raise learning behavior \emph{only}
through psychological safety, and that learning behavior in turn raises
team performance. All claims are proved from the explicitly stated
behavioral axioms.
\end{abstract}

\section{Formalization}

\subsection{Units and levels}
We consider a population $\mathcal{T}$ of ongoing work teams $t$, each
embedded in an organization and persisting over time, with member set
$M_t$, $|M_t|=n_t\ge 2$. Quantities indexed by $i$ range over individual
members $i\in M_t$; quantities indexed by $t$ are \emph{team-level}
(either properties of the team as a whole, e.g.\ leader behavior, or
aggregates of member-level quantities). This separation of levels is the
formal counterpart of the requirement to identify individual-, team-, and
organizational-level antecedents.

\subsection{The dependent construct: learning behavior}
\begin{definition}[Learning behavior]\label{def:LB}
For team $t$, let $B_t=(B_t^{1},\dots,B_t^{5})$ record the team's
\emph{learning behavior}, where the five coordinates are the (suitably
scaled, e.g.\ frequency-per-unit-time) intensities of the five activities
named in the problem:
\[
B_t^{1}=\text{asking questions},\quad
B_t^{2}=\text{seeking feedback},\quad
B_t^{3}=\text{discussing errors/unexpected outcomes},
\]
\[
B_t^{4}=\text{reflecting on past performance},\quad
B_t^{5}=\text{experimenting with new actions/methods}.
\]
Define the scalar \emph{learning-behavior level}
\[
L_t \;=\; \sum_{k=1}^{5} w_k\, B_t^{k},\qquad w_k>0,\ \sum_k w_k=1,
\]
a positively weighted index of the five activities. All five activities
are observable, so $L_t$ is an observable random variable.
\end{definition}

\begin{remark}
The defining feature shared by $B_t^1,\dots,B_t^5$ is that each requires a
member to expose a potential deficiency or to deviate from the routine in
front of others: a question reveals ignorance, feedback-seeking reveals
uncertainty about one's own performance, discussing an error reveals a
mistake, reflection surfaces collective shortcomings, and experimentation
risks visible failure. Hence every coordinate of $B_t$ is an
\emph{interpersonally risky} act. This single structural fact---proved as
Lemma~\ref{lem:risk} below from the definitions---is what makes an
interpersonal-climate variable the natural mediator.
\end{remark}

\subsection{The mediator: team psychological safety}
\begin{definition}[Psychological safety]\label{def:PS}
For member $i$ in team $t$ let $\psi_{t,i}\in\mathbb{R}$ denote $i$'s
belief about the consequences of taking interpersonal risks in the team,
measured so that larger values mean ``risk taking is safer'' (smaller
expected interpersonal cost). The team-level \emph{psychological safety}
is the shared component
\[
\Psi_t \;=\; \frac{1}{n_t}\sum_{i\in M_t}\psi_{t,i},
\]
and we say the belief is \emph{shared} when the within-team dispersion
$\operatorname{Var}_{i\in M_t}(\psi_{t,i})$ is small (formalized in
Assumption~\ref{ass:share}). $\Psi_t$ is the ``shared interpersonal
climate variable'' of the problem.
\end{definition}

\subsection{The exogenous contextual supports}
\begin{definition}[Contextual supports]\label{def:X}
Let $X_t=(S_t,C_t,K_t)$ be team-level antecedents:
$S_t$ supportive leader behavior, $C_t$ organizational context support
(resources, information, rewards aligned with learning), and $K_t$
coaching. These are the ``contextual supports'' of the problem and live at
the team/organizational level. Write $X_t\in\mathbb{R}^3$ and let
$\beta=(\beta_S,\beta_C,\beta_K)$ be a coefficient vector.
\end{definition}

\subsection{Individual-level antecedents and performance}
\begin{definition}\label{def:Z}
Let $Z_{t,i}$ collect individual-level antecedents of member $i$
(e.g.\ tenure, self-efficacy, status). Let $P_t$ denote a team
performance outcome (e.g.\ supervisor-rated effectiveness), scaled so that
larger is better.
\end{definition}

\section{The structural model}

We posit a recursive (acyclic) linear structural system. All disturbance
terms $\varepsilon_\bullet$ are mean-zero and, by the recursive ordering
$X_t\to \psi_{t,i}\to L_t\to P_t$, are uncorrelated with the regressors
that precede them in the ordering (Assumption~\ref{ass:exog}).

\begin{assumption}[Risk-cost formation of beliefs]\label{ass:safety}
For each member,
\begin{equation}\label{eq:psi}
\psi_{t,i} \;=\; \alpha_0 + \beta^{\top} X_t + \gamma^{\top} Z_{t,i}
             + \varepsilon^{\psi}_{t,i},
\end{equation}
with $\beta_S,\beta_C,\beta_K>0$. That is, each contextual support
raises members' beliefs that interpersonal risk taking is safe.
\end{assumption}

\begin{assumption}[Shared climate]\label{ass:share}
Members of a team share a common response to team-level inputs: the slopes
$\beta,\gamma$ in \eqref{eq:psi} are common across $i\in M_t$, and the
idiosyncratic shocks $\varepsilon^{\psi}_{t,i}$ are i.i.d.\ across $i$
within a team, independent of $(X_t,\bar Z_t)$, with mean $0$ and variance
$\sigma^2_\psi$. Write the team's \emph{systematic component}
\[
\tau_t \;=\; \alpha_0+\beta^\top X_t+\gamma^\top\bar Z_t,
\qquad
\bar Z_t=\tfrac1{n_t}\textstyle\sum_{i\in M_t} Z_{t,i},
\]
and let $\tau^2=\operatorname{Var}_{t}(\tau_t)$ denote its variance across
teams in the population $\mathcal T$. Two consequences follow.
\emph{(a) Within-team dispersion is bounded:} since members of a team share
the common $\tau_t$, the only source of within-team variation in
$\psi_{t,i}=\tau_t+\gamma^\top(Z_{t,i}-\bar Z_t)+\varepsilon^{\psi}_{t,i}$
is idiosyncratic, so $\Psi_t$ is a meaningful shared-climate aggregate
whenever $\tau^2$ is large relative to $\sigma^2_\psi$.
\emph{(b) Aggregation is reliable:} averaging \eqref{eq:psi} over $i$,
$\Psi_t=\tau_t+\tfrac1{n_t}\sum_i\varepsilon^{\psi}_{t,i}$, so the
reliability of the team-mean climate (the squared correlation between
$\Psi_t$ and its systematic target $\tau_t$) is
\[
\rho \;=\;
\frac{\operatorname{Var}_t(\tau_t)}
     {\operatorname{Var}_t(\tau_t)+\sigma^2_\psi/n_t}
\;=\;\frac{\tau^2}{\tau^2+\sigma^2_\psi/n_t}\;\longrightarrow\;1
\]
as $n_t\to\infty$ or as the systematic between-team variance $\tau^2$
dominates $\sigma^2_\psi$. Thus $\Psi_t$ is a well-defined shared
team-level variable and not an artifact of averaging noise.
\end{assumption}

\begin{assumption}[Behavior responds to climate, not directly to supports]%
\label{ass:behavior}
Aggregated learning behavior is generated by
\begin{equation}\label{eq:L}
L_t \;=\; \lambda_0 + \delta\,\Psi_t + \theta^{\top}\bar Z_t
          + \varepsilon^{L}_t,
\end{equation}
with $\delta>0$, and \emph{no} direct path from $X_t$ to $L_t$ once
$\Psi_t$ is included: $\partial L_t/\partial X_t \,|\, \Psi_t = 0$.
\end{assumption}

\begin{assumption}[Performance responds to learning]\label{ass:perf}
\begin{equation}\label{eq:P}
P_t \;=\; \pi_0 + \mu\, L_t + \nu^{\top} W_t + \varepsilon^{P}_t,
\qquad \mu>0,
\end{equation}
where $W_t$ are performance covariates excluded from \eqref{eq:L}.
\end{assumption}

\begin{assumption}[Recursive exogeneity]\label{ass:exog}
$\mathbb{E}[\varepsilon^{\psi}_{t,i}\mid X_t,Z_{t,i}]=0$,
$\mathbb{E}[\varepsilon^{L}_t\mid \Psi_t,\bar Z_t]=0$, and
$\mathbb{E}[\varepsilon^{P}_t\mid L_t,W_t]=0$; moreover
$\varepsilon^L_t$ is uncorrelated with $X_t$ given $\Psi_t,\bar Z_t$.
\end{assumption}

These four assumptions are the behavioral axioms; everything below is
derived from them.

\section{Justification of the structural links (mechanisms)}

\begin{lemma}[Every learning activity is interpersonally risky]\label{lem:risk}
Each coordinate $B_t^k$, $k=1,\dots,5$, of Definition~\ref{def:LB} is an
act whose execution by a member can, in the absence of safety, incur an
interpersonal cost (being seen as ignorant, incompetent, negative, or
disruptive). Hence the disposition to perform $B_t^k$ is decreasing in the
\emph{perceived} interpersonal cost of risk taking and increasing in
$\psi_{t,i}$.
\end{lemma}

\begin{proof}
We argue activity by activity. (1) Asking a question makes one's lack of
knowledge observable to others; the possible cost is the image of
ignorance. (2) Seeking feedback exposes uncertainty about one's own
performance and invites evaluation; the cost is a negative appraisal. (3)
Openly discussing an error or unexpected outcome makes a failure
attributable to identifiable members; the cost is blame. (4) Reflecting on
past performance surfaces collective and individual shortcomings for
public examination; the cost is criticism. (5) Experimenting with a new
method departs from the sanctioned routine and may fail visibly; the cost
is the image of incompetence or the charge of being disruptive. In each
case the act has the form ``expose a potential deficiency / deviation
before others,'' which is the definition of an interpersonal risk. A
rational member weighs the expected interpersonal cost against the benefit;
since $\psi_{t,i}$ measures (inversely) that expected cost, the propensity
to act is monotone increasing in $\psi_{t,i}$. As $L_t$ in
Definition~\ref{def:LB} is a positive combination of the five activities,
the team propensity for $L_t$ inherits this monotonicity. This is exactly
the qualitative content encoded by $\delta>0$ in \eqref{eq:L}.
\end{proof}

\begin{lemma}[Contextual supports raise safety]\label{lem:supports}
Supportive leadership, context support, and coaching each lower the
expected interpersonal cost of risk taking, hence raise $\psi_{t,i}$:
the signs $\beta_S,\beta_C,\beta_K>0$ in \eqref{eq:psi} are justified.
\end{lemma}

\begin{proof}
(i) Supportive leader behavior ($S_t$): a leader who is accessible,
invites input, and does not punish messengers reduces the probability that
a question, an admitted error, or a failed experiment is met with
sanction; the leader's behavior is the most salient signal members use to
estimate interpersonal cost, so increasing $S_t$ decreases that estimate,
i.e.\ raises $\psi_{t,i}$ ($\beta_S>0$). (ii) Context support ($C_t$):
resources, information access, and rewards aligned with learning make risk
taking less likely to be futile or penalized and more likely to be
rewarded, lowering its net expected cost ($\beta_C>0$). (iii) Coaching
($K_t$): coaching frames difficulties as shared learning problems, models
inquiry, and provides developmental rather than evaluative feedback, again
lowering the expected interpersonal cost ($\beta_K>0$). Each mechanism acts
on the \emph{belief about consequences of risk taking}, which is precisely
$\psi_{t,i}$; none of them, by Assumption~\ref{ass:behavior}, acts on
$L_t$ except through that belief. The shared-slope condition
(Assumption~\ref{ass:share}) lets these team-level inputs move the team
\emph{mean} $\Psi_t$, giving $\partial \Psi_t/\partial X_t=\beta$.
\end{proof}

\begin{lemma}[Learning behavior raises performance]\label{lem:perf}
$\mu>0$ in \eqref{eq:P}: greater learning behavior improves team
performance.
\end{lemma}

\begin{proof}
Ongoing teams in real organizations face changing tasks, technologies, and
demands. The five activities each generate or apply knowledge that adjusts
the team's process to current conditions: questions and feedback import
missing information; discussing errors and reflecting convert past
outcomes into corrected routines; experimentation discovers superior
methods. Each adjustment weakly improves the match between the team's
process and its task environment, and at least one is strictly improving
whenever the environment is non-stationary; aggregating over the
activities gives a positive marginal effect of $L_t$ on the performance
outcome $P_t$, i.e.\ $\mu>0$. The covariates $W_t$ absorb performance
determinants unrelated to learning so that $\mu$ isolates the
learning channel (Assumption~\ref{ass:exog}).
\end{proof}

\section{The mediation theorem}

We use the standard definition of (complete) mediation in a recursive
linear system.

\begin{definition}[Direct, indirect, total effects]\label{def:effects}
In the system \eqref{eq:psi}--\eqref{eq:L}, holding individual antecedents
fixed at their team means, the effect of $X_t$ on $L_t$ decomposes as
\[
\underbrace{\frac{dL_t}{dX_t}}_{\text{total}}
=\underbrace{\frac{\partial L_t}{\partial X_t}\Big|_{\Psi_t}}_{\text{direct}}
+\underbrace{\frac{\partial L_t}{\partial \Psi_t}\,\frac{\partial \Psi_t}{\partial X_t}}_{\text{indirect}} .
\]
$\Psi_t$ \emph{completely mediates} the effect of $X_t$ on $L_t$ if the
direct effect is $0$ and the indirect effect is nonzero.
\end{definition}

\begin{theorem}[Psychological safety completely mediates the support--learning link]%
\label{thm:main}
Under Assumptions~\ref{ass:safety}--\ref{ass:exog}, the effect of the
contextual supports $X_t=(S_t,C_t,K_t)$ on learning behavior $L_t$ is
\begin{equation}\label{eq:decomp}
\frac{dL_t}{dX_t} \;=\; \delta\,\beta \;\neq\; 0,
\qquad
\frac{\partial L_t}{\partial X_t}\Big|_{\Psi_t}=0,
\end{equation}
so $\Psi_t$ completely mediates the relationship. Each component effect is
strictly positive:
\[
\frac{dL_t}{dS_t}=\delta\beta_S>0,\quad
\frac{dL_t}{dC_t}=\delta\beta_C>0,\quad
\frac{dL_t}{dK_t}=\delta\beta_K>0.
\]
Moreover the supports raise performance, again entirely through the chain
$X_t\to\Psi_t\to L_t\to P_t$:
\begin{equation}\label{eq:perfchain}
\frac{dP_t}{dX_t} \;=\; \mu\,\delta\,\beta,\qquad
\frac{dP_t}{dS_t}=\mu\delta\beta_S>0,\ \text{etc.}
\end{equation}
\end{theorem}

\begin{proof}
\emph{Step 1 (mediator is driven by supports).}
Average \eqref{eq:psi} over $i\in M_t$. By Assumption~\ref{ass:share} the
slopes $\beta,\gamma$ are common, so
\[
\Psi_t=\frac1{n_t}\sum_i\psi_{t,i}
=\alpha_0+\beta^\top X_t+\gamma^\top \bar Z_t
+\frac1{n_t}\sum_i\varepsilon^\psi_{t,i}.
\]
Hence $\partial\Psi_t/\partial X_t=\beta=(\beta_S,\beta_C,\beta_K)$, which
is strictly positive in each coordinate by Lemma~\ref{lem:supports}.

\emph{Step 2 (no direct path).}
By Assumption~\ref{ass:behavior}, $L_t$ depends on $X_t$ only through
$\Psi_t$ once $\Psi_t$ and $\bar Z_t$ are conditioned on:
$\partial L_t/\partial X_t|_{\Psi_t}=0$. This is the second equality in
\eqref{eq:decomp}.

\emph{Step 3 (indirect effect).}
From \eqref{eq:L}, $\partial L_t/\partial\Psi_t=\delta>0$ by
Lemma~\ref{lem:risk}. Substituting $\Psi_t$ from Step~1 into \eqref{eq:L},
\begin{equation}\label{eq:reduced}
L_t=\big(\lambda_0+\delta\alpha_0\big)
+\delta\beta^\top X_t
+\big(\theta+\delta\gamma\big)^\top\bar Z_t
+\delta\,\tfrac1{n_t}\textstyle\sum_i\varepsilon^\psi_{t,i}
+\varepsilon^L_t .
\end{equation}
Differentiating, $dL_t/dX_t=\delta\beta$. Combining with Steps~1--2 gives
\eqref{eq:decomp}: total effect $=$ indirect effect $=\delta\beta$, direct
effect $=0$. Each coordinate equals $\delta\beta_S,\delta\beta_C,
\delta\beta_K$, all strictly positive since $\delta>0$ and
$\beta_S,\beta_C,\beta_K>0$. The indirect effect is nonzero, so by
Definition~\ref{def:effects} $\Psi_t$ is a complete mediator.

\emph{Step 4 (carry through to performance).}
Substitute \eqref{eq:reduced} into \eqref{eq:P}:
\[
P_t=\big(\pi_0+\mu\lambda_0+\mu\delta\alpha_0\big)
+\mu\delta\beta^\top X_t
+\mu(\theta+\delta\gamma)^\top\bar Z_t
+\nu^\top W_t+ u_t,
\]
where $u_t=\mu\delta\frac1{n_t}\sum_i\varepsilon^\psi_{t,i}+\mu\varepsilon^L_t+\varepsilon^P_t$.
Thus $dP_t/dX_t=\mu\delta\beta$, strictly positive in each coordinate
because $\mu>0$ (Lemma~\ref{lem:perf}), $\delta>0$, $\beta\gg 0$. This is
\eqref{eq:perfchain}, and it shows the entire effect of contextual
supports on performance flows through the chain
$X_t\to\Psi_t\to L_t\to P_t$.

\emph{Step 5 (identification / unbiasedness).}
By Assumption~\ref{ass:exog} the disturbance in each structural equation is
mean-independent of that equation's regressors, so the population
regressions of $\Psi_t$ on $(X_t,\bar Z_t)$, of $L_t$ on
$(\Psi_t,\bar Z_t)$, and of $P_t$ on $(L_t,W_t)$ recover
$\beta,\delta,\mu$ respectively without bias; in particular the regression
of $L_t$ on $(X_t,\Psi_t,\bar Z_t)$ returns coefficient $0$ on $X_t$ and
$\delta$ on $\Psi_t$, confirming complete mediation is testable and
identified, not merely assumed.
\end{proof}

\begin{corollary}[Partial mediation as the general case]\label{cor:partial}
If Assumption~\ref{ass:behavior} is relaxed to admit a direct term
$\zeta^\top X_t$ in \eqref{eq:L}, then
$dL_t/dX_t=\zeta+\delta\beta$ with indirect part $\delta\beta$ and direct
part $\zeta$; $\Psi_t$ is a \emph{partial} mediator iff $\delta\beta\neq0$
and $\zeta\neq0$, and the complete-mediation conclusion of
Theorem~\ref{thm:main} is the special case $\zeta=0$. Thus the model
nests both regimes, with the sign and magnitude of the mediated portion
$\delta\beta$ always positive under
Lemmas~\ref{lem:risk}--\ref{lem:supports}.
\end{corollary}

\begin{proof}
Immediate: replace $0$ by $\zeta^\top X_t$ in
Assumption~\ref{ass:behavior} and repeat Steps~1--3; the reduced form
\eqref{eq:reduced} gains a term $\zeta^\top X_t$, giving
$dL_t/dX_t=\zeta+\delta\beta$. The classification follows from
Definition~\ref{def:effects}.
\end{proof}

\section{Multilevel summary of antecedents and mechanisms}

The model answers each part of the problem explicitly.

\begin{itemize}
\item \textbf{Individual level.} The antecedents $Z_{t,i}$ enter
\eqref{eq:psi} and, via their team mean $\bar Z_t$, \eqref{eq:L}.
Mechanism: a member's tenure, self-efficacy, and status calibrate the
\emph{personal} expected cost of interpersonal risk
($\gamma^\top Z_{t,i}$ in \eqref{eq:psi}) and the personal capacity to act
($\theta^\top\bar Z_t$ in \eqref{eq:L}). They modulate but do not
substitute for the shared climate.

\item \textbf{Team level.} Supportive leadership $S_t$ and coaching $K_t$
are team-level inputs whose mechanism (Lemma~\ref{lem:supports}) is to
lower the shared expected cost of risk taking, raising $\Psi_t$; and the
shared climate $\Psi_t$ itself is the team-level mediator whose mechanism
(Lemma~\ref{lem:risk}) is to make the interpersonally risky learning
activities $B_t^1,\dots,B_t^5$ enactable.

\item \textbf{Organizational level.} Context support $C_t$
(resources, information, learning-aligned rewards) is the
organizational-level input; its mechanism (Lemma~\ref{lem:supports}) is to
reduce the futility and penalty of risk taking, again operating through
$\Psi_t$.

\item \textbf{Mediation.} Theorem~\ref{thm:main} proves that
$\Psi_t$ completely mediates contextual supports $\to$ learning behavior
(direct effect $0$, indirect effect $\delta\beta\gg0$), with the general
partial-mediation regime recovered in Corollary~\ref{cor:partial}.

\item \textbf{Performance.} Equation~\eqref{eq:P} and
Lemma~\ref{lem:perf} establish $\Psi_t\to L_t\to P_t$ with total support
effect $\mu\delta\beta\gg0$ \eqref{eq:perfchain}: the climate variable
improves performance precisely by enabling learning behavior, not
independently of it.
\end{itemize}

This completes the construction and justification of the model.
\end{document}
